% Partie : intuitionistic-logic
% Chapitre : semantics
% Section : topological-semantics

\documentclass[../../../include/open-logic-section]{subfiles}

\begin{document}

\olfileid[fr]{int}{sem}{top}

\olsection{Sémantique topologique}

Une autre façon de donner une sémantique à la logique intuitionniste
repose sur la notion mathématique de topologie.

\begin{defn}
  Soit $X$ un ensemble. Une \emph{topologie sur~$X$} est un ensemble
  $\Top{O} \subseteq \Pow{X}$ qui satisfait les propriétés ci-dessous.
  Les !!{element}s de~$\Top{O}$ sont appelés les \emph{ouverts} de
  cette topologie. L'ensemble $X$ muni de $\Top{O}$ est appelé un
  \emph{espace topologique}.
  \begin{enumerate}
  \item L'ensemble vide et l'espace entier sont ouverts : $\emptyset$, $X \in
    \Top{O}$.
  \item Les ouverts sont stables par intersections finies : si $U$, $V \in
    \Top{O}$, alors $U \cap V \in \Top{O}$.
  \item Les ouverts sont stables par réunions arbitraires : si $U_i \in
    \Top{O}$ pour tout $i \in I$, alors $\bigcup \Setabs{U_i}{i \in
      I} \in \Top{O}$.
  \end{enumerate}
\end{defn}

Par abus de notation, nous pouvons écrire $X$ pour désigner l'espace
topologique lorsque la collection de ses ouverts ressort du contexte.
Il reste que $X$ désigne l'ensemble sous-jacent, et que la topologie
est la collection d'ouverts dont il est muni.

\begin{defn}
  Un \emph{modèle topologique} de la logique propositionnelle
  intuitionniste est un triplet $\mModel{X} = \tuple{X, \Top{O}, V}$ où
  $\Top{O}$ est une topologie sur~$X$ et $V$ est une fonction qui
  associe à chaque variable propositionnelle un ouvert de
  $\Top{O}$.

  Étant donné un modèle topologique~$\mModel{X}$, nous définissons
  $\Prop{X}{!A}$ par induction comme suit :
  \begin{enumerate}
  \item $\Prop{X}{\lfalse} = \emptyset$
  \item $\Prop{X}{p} = V(p)$
  \item $\Prop{X}{!A \land !B} = \Prop{X}{!A} \cap \Prop{X}{!B}$
  \item $\Prop{X}{!A \lor !B} = \Prop{X}{!A} \cup \Prop{X}{!B}$
  \item $\Prop{X}{!A \lif !B} = \Interior{(X \setminus \Prop{X}{!A}) \cup
    \Prop{X}{!B}}$
  \end{enumerate}
  Ici, $\Interior{V}$ associe à tout ensemble $V \subseteq X$ son
  \emph{intérieur}, c'est-à-dire la réunion de tous les ouverts qu'il
  contient. Autrement dit,
  \[
  \Interior{V} = \bigcup \Setabs{U}{U \subseteq V \text{ et } U \in \Top{O}}.
  \]
\end{defn}

L'intérieur de tout ensemble est ouvert, puisqu'il est la réunion
d'ouverts. Ainsi, $\Prop{X}{!A}$ est toujours un ouvert.

Cette sémantique est abstraite, mais on peut en éclairer l'idée.
Supposons que les !!{element}s, ou « points », de $X$ soient des lieux
où l'on peut évaluer les énoncés. L'ensemble des points où $!A$ est
vraie est la proposition exprimée par~$!A$. Tous les ensembles de
points ne peuvent pas jouer ce rôle : seuls les !!{element}s de
$\Top{O}$ le peuvent. On a $!A \Entails !B$ si et seulement si $!B$
est vraie en tout point où~$!A$ est vraie, c'est-à-dire si
$\Prop{X}{!A} \subseteq \Prop{X}{!B}$ dans chaque modèle sur un
espace~$X$. L'énoncé absurde~$\lfalse$ n'est vrai en aucun point ;
donc $\Prop{X}{\lfalse} = \emptyset$.

Quel rapport faut-il établir entre les propositions exprimées par
$!B \land !C$, $!B \lor !C$ et $!B \lif !C$, et celles exprimées par
$!B$ et~$!C$, pour respecter les lois intuitionnistes valides ? Le
critère recherché est que $!A \Proves !B$ si et seulement si
$\Prop{X}{!A} \subseteq \Prop{X}{!B}$ dans tous les modèles
topologiques. Nous voulons $\lfalse \Proves !A$ pour toute $!A$,
ce qui est vérifié puisque $\emptyset \subseteq U$ pour tout $U$.
Comme $!B \land !C \Proves !B$, il faut
$\Prop{X}{!B \land !C} \subseteq \Prop{X}{!B}$ et, de même,
$\Prop{X}{!B \land !C} \subseteq \Prop{X}{!C}$. Le plus grand
ensemble satisfaisant $W \subseteq U$ et $W \subseteq V$ est
$U \cap V$. Réciproquement, $!B \Proves !B \lor !C$ et
$!C \Proves !B \lor !C$ ; il faut donc
$\Prop{X}{!B} \subseteq \Prop{X}{!B \lor !C}$ et
$\Prop{X}{!C} \subseteq \Prop{X}{!B \lor !C}$. Le plus petit
ensemble~$W$ tel que $U \subseteq W$ et $V \subseteq W$ est
$U \cup V$.

La définition pour $\lif$ est plus délicate : $!A \lif !B$ exprime la
proposition la plus faible qui, jointe à $!A$, entraîne $!B$. Le fait
que $!A \lif !B$ joint à $!A$ entraîne~$!B$ découle de
$(!A \lif !B) \land !A \Proves !B $. Ainsi, $\Prop{X}{!A \lif !B}$
doit être le plus grand ouvert dont l'intersection avec
$\Prop{X}{!A}$ est incluse dans $\Prop{X}{!B}$. En effet, pour
des ouverts $U$ et $V$, la condition $W\cap U\subseteq V$
équivaut à $W\subseteq (X\setminus U)\cup V$. Si $W$ est ouvert,
cette dernière condition équivaut à
$W\subseteq \Interior{(X\setminus U)\cup V}$. L'intérieur est
donc le plus grand ouvert qui convienne.

\begin{ex}
  Prenons $X=\Real$ muni de sa topologie usuelle et
  une valuation telle que $V(p)=(0,+\infty)$. Comme $\lnot p$
  abrège $p\lif\lfalse$, la clause d'implication donne
  \[
    \Prop{X}{\lnot p} = \Interior{(-\infty,0]} = (-\infty,0).
  \]
  Par conséquent, $\Prop{X}{p\lor\lnot p} = \Real\setminus\{0\}$.
  Au point~$0$, ni $p$ ni $\lnot p$ ne sont vraies. Ce modèle
  montre que le tiers exclu n'est pas valide dans la sémantique
  topologique.
\end{ex}

\end{document}
